\documentclass[11pt]{amsart}
\usepackage[margin=1.15in]{geometry}
\usepackage{amsmath,amssymb,amsthm,mathtools}
\usepackage[T1]{fontenc}
\usepackage{lmodern}
\usepackage{microtype}
\usepackage{enumitem}
\usepackage{alltt}
\usepackage{xcolor}
\usepackage[colorlinks=true,linkcolor=blue!50!black,citecolor=blue!50!black,urlcolor=blue!50!black]{hyperref}

\newtheorem{theorem}{Theorem}[section]
\newtheorem{lemma}[theorem]{Lemma}
\newtheorem{proposition}[theorem]{Proposition}
\newtheorem{corollary}[theorem]{Corollary}
\newtheorem{maintheorem}{Theorem}
\renewcommand{\themaintheorem}{\Alph{maintheorem}}
\newtheorem{mainproposition}[maintheorem]{Proposition}
\theoremstyle{definition}
\newtheorem{definition}[theorem]{Definition}
\newtheorem{question}{Question}
\renewcommand{\thequestion}{\Alph{question}}
\theoremstyle{remark}
\newtheorem{remark}[theorem]{Remark}

\newcommand{\R}{\mathbb{R}}
\newcommand{\N}{\mathbb{N}}
\newcommand{\dimH}{\dim_{\mathrm H}}
\newcommand{\Lm}{\mathcal{L}}
\newcommand{\Dy}{\mathcal{D}}


\title[Erd\H{o}s similarity for distance sets]{Erd\H{o}s similarity for distance sets:\\
finite patterns, a pinned obstruction, and an open window}
\author{Harald Schilly}
\address{Independent researcher, Vienna, Austria}
\email{harald.schilly@gmail.com}
\date{October 2026}

\begin{document}

\begin{abstract}
For a compact set $E\subset\R^d$, $d\ge2$, let $\Delta(E)=\{|x-y|:x,y\in E\}$ be its distance set.
We ask whether distance sets of sets of large Hausdorff dimension must contain a nontrivial
affine copy $x+s\{2^{-n}:n\ge1\}$ of the dyadic sequence. This is a distance-set version of
the Erd\H{o}s similarity problem. We combine two recent results announced by OpenAI: the Falconer
distance theorem in all dimensions, and the dyadic case of the Erd\H{o}s similarity conjecture.
Assuming these, we show the following. (A) If $\dimH E>d/2$, then $\Delta(E)$ contains affine
copies of every finite set, for a positive-measure family of translations at every small scale.
(B) There is a compact set $K\subset\R^d$ of positive volume, hence of full dimension, whose
distance set from the origin contains no affine copy of the dyadic sequence. The full distance
set $\Delta(K)$ does contain one, and the origin is the only point of $\R^d$ ($d\ge2$) from
which the obstruction occurs. Since the Mattila--Sj\"olin theorem settles the question above dimension
$(d+1)/2$, it remains open exactly in the window $d/2<\dimH E\le(d+1)/2$. A negative answer
there would produce distance sets of positive measure with empty interior. Both theorems are
formalized in Lean~4 with Mathlib, with the two inputs as explicit hypotheses. We also
re-checked OpenAI's formal proofs of both inputs with the Lean kernel.
\end{abstract}

\subjclass[2020]{28A80 (primary); 28A75, 42B10, 03B35}
\keywords{Erd\H{o}s similarity problem, distance sets, Falconer conjecture, Hausdorff dimension,
formal verification, Lean}

\maketitle

\section{Introduction}

\subsection{Two problems about large sets}
A set $P\subset\R$ is \emph{measure universal} if every measurable $S\subset\R$ of positive
Lebesgue measure contains a nontrivial affine copy $x+sP$ ($x\in\R$, $s\neq0$). Steinhaus'
theorem~\cite{Steinhaus1920} implies that every finite set is measure universal. Erd\H{o}s
conjectured~\cite{Erdos1974} that no infinite set is. The conjecture is known for slowly
decaying sequences~\cite{Falconer1984} and for many structured sets~\cite{Bourgain1987,Kolountzakis1997}.
The geometric sequences $\{q^n\}$, and above all the dyadic sequence
\[
  \Dy=\{2^{-n}:n\ge1\},
\]
were long described as the main bottleneck~\cite{JungLaiMooroogen2024}. In September and October
2026, OpenAI released a collection of manuscripts produced by an internal AI model. It
includes a proof that $\Dy$ is not measure universal~\cite{OAIDyadic} and an extension to every
ratio $q\in(0,1)$~\cite{OAIGeometric}.

Falconer's distance problem~\cite{Falconer1985} asks for the smallest $\alpha$ such that every
compact $E\subset\R^d$ with $\dimH E>\alpha$ has a distance set
\[
  \Delta(E)=\{|x-y|:x,y\in E\}
\]
of positive Lebesgue measure. Falconer showed that $d/2$ would be optimal and conjectured it.
Partial progress includes~\cite{Wolff1999,Erdogan2005,GIOW2020}. The same OpenAI release contains a
proof of the full conjecture in every dimension $d\ge2$~\cite{OAIFalconer}.
Mattila and Sj\"olin~\cite{MattilaSjolin1999} showed that $\dimH E>(d+1)/2$ even forces
$\Delta(E)$ to contain an interval.

\subsection{The question}
The Falconer theorem says that $\Delta(E)$ is large in measure, while the dyadic theorem says
that being large in measure is not enough to contain $\Dy$. Distance sets, however, have
additional multi-scale structure: $\Delta(E)\supseteq\Delta(E\cap B)$ for every ball $B$, so
$\Delta(E)\cap[0,r)$ has positive measure for every $r>0$ as soon as the dimension condition
holds locally. This motivates:

\begin{question}[Erd\H{o}s similarity for distance sets]\label{q:main}
Let $d\ge2$ and let $E\subset\R^d$ be compact with $\dimH E>d/2$. Must $\Delta(E)$ contain a
nontrivial affine copy of $\Dy$?
\end{question}

We could not find this question, or any study of Erd\H{o}s-type patterns inside distance sets,
in the literature; it is not discussed in the recent survey~\cite{JungLaiMooroogen2024}.

\subsection{Results}
Throughout, Theorems~A and~B are \emph{conditional} on the two inputs stated precisely in
Section~\ref{sec:inputs}. Theorem~\ref{thm:inputFalconer} is the Falconer theorem of~\cite{OAIFalconer}, and
Theorem~\ref{thm:inputDyadic} is the dyadic avoidance theorem of~\cite{OAIDyadic}.

\begin{maintheorem}[Finite patterns]\label{thm:A}
Assume Theorem~\ref{thm:inputFalconer}. Let $d\ge2$, let $E\subset\R^d$ be compact with $\dimH E>d/2$, and let
$F\subset\R$ be finite. There is $\delta>0$ such that for every $|s|<\delta$,
\[
  \Lm^1\bigl(\{x\in\R : x+sF\subseteq\Delta(E)\}\bigr)>0 .
\]
In particular $\Delta(E)$ contains affine copies $x+sF$ with $0<s<\delta'$ and with
$-\delta'<s<0$, for every $\delta'>0$.
\end{maintheorem}

\begin{maintheorem}[A pinned obstruction]\label{thm:B}
Assume Theorem~\ref{thm:inputDyadic}. Let $d\ge1$ and $\eta\in(0,1)$. There is a compact set $K$ in the closed
unit ball of $\R^d$ with $0\in K$ and
\[
  \Lm^d(K)\ \ge\ (1-\eta)^d\,\Lm^d(B(0,1))>0,\qquad\text{hence }\dimH K=d,
\]
such that the pinned distance set $\Delta_0(K)=\{|y|:y\in K\}$ contains no nontrivial affine
copy of $\Dy$, whereas the full distance set $\Delta(K)$ contains one.
If $d\ge2$, then for every $p\in\R^d\setminus\{0\}$ the pinned distance set
$\Delta_p(K)=\{|p-y|:y\in K\}$ contains a nondegenerate interval, hence an affine copy of $\Dy$.
\end{maintheorem}

So pinned distance sets are not universal hosts for $\Dy$, even from a pin inside a set of
full dimension and nearly full volume in the ball. The obstruction, however, sits at one single
point. Theorems~\ref{thm:A} and~\ref{thm:B} together with the Mattila--Sj\"olin theorem show
where Question~\ref{q:main} is undecided:

\begin{mainproposition}\label{prop:C}
Let $d\ge2$ and let $E\subset\R^d$ be compact. If $\dimH E>(d+1)/2$, or if $\Lm^d(E)>0$, then
$\Delta(E)$ contains an affine copy $a+s\,P$ of every bounded set $P\subset[0,\infty)$; in
particular of $\Dy$. If Question~\ref{q:main} has a negative answer for some $E$ with
$d/2<\dimH E\le(d+1)/2$, then $\Delta(E)$ has positive measure but empty interior.
\end{mainproposition}

We are not aware of a general result that guarantees nonempty interior of $\Delta(E)$ below the
Mattila--Sj\"olin threshold $(d+1)/2$ (see~\cite{MattilaBook2015} for background). Under that
reading, a negative answer to Question~\ref{q:main} in the window would also give new
examples of distance sets with positive measure and empty interior. Section~\ref{sec:open}
lists further questions.

\subsection{Formalization and status of the inputs}
Theorems~\ref{thm:A} and~\ref{thm:B} (the latter without the explicit constant $(1-\eta)^d$
and without the last sentence; see Section~\ref{sec:lean}) are formalized in Lean~4~\cite{Lean4}
with Mathlib~\cite{Mathlib2020}. The two inputs enter as explicit hypotheses whose Lean types are
copied verbatim from OpenAI's formal statements, and the formal proofs use no axioms beyond
Lean's standard ones. The two inputs come from an AI-generated, unrefereed manuscript collection
whose authors state that some results ``could have issues''. OpenAI also provides Lean proofs of both
inputs. We rebuilt the relevant part of that library ourselves and checked with the Lean kernel
that these proofs establish exactly our hypotheses, using only the standard axioms
(Section~\ref{sec:lean}). This makes the dependence on the inputs as transparent as we can.
It does not replace independent mathematical review of~\cite{OAIFalconer,OAIDyadic}.

\section{Notation and the two inputs}\label{sec:inputs}
$\Lm^d$ is Lebesgue measure on $\R^d$, $|\cdot|$ the Euclidean norm, $B(0,1)$ the open unit
ball, and $\dimH$ Hausdorff dimension. For $E\subset\R^d$ and $p\in\R^d$ we write
$\Delta(E)=\{|x-y|:x,y\in E\}$ and $\Delta_p(E)=\{|p-y|:y\in E\}$. An \emph{affine copy} of
$P\subset\R$ is a set $x+sP$ with $x\in\R$ and $s\ne0$.

\begin{theorem}[Falconer distance theorem~{\cite{OAIFalconer}}]\label{thm:inputFalconer}
Let $d\ge2$ and let $E\subset\R^d$ be compact with $\dimH E>d/2$. Then $\Lm^1(\Delta(E))>0$.
\end{theorem}

\begin{theorem}[Dyadic avoidance~{\cite{OAIDyadic}}]\label{thm:inputDyadic}
For every $\eta\in(0,1)$ there is a compact set $A\subseteq[0,1]$ with $\Lm^1(A)>1-\eta$ such that
for every $x\in\R$ and every $s\neq0$ there is $n\ge1$ with $x+s2^{-n}\notin A$.
\end{theorem}

\section{Finite patterns: proof of Theorem~\ref{thm:A}}

The following is the standard Steinhaus-type argument (compare the introduction
of~\cite{OAIDyadic}). We include it because Theorem~\ref{thm:A} needs the version with a
positive-measure set of translations at every small scale.

\begin{lemma}\label{lem:steinhaus}
Let $S\subset\R$ be measurable with $0<\Lm^1(S)<\infty$ and let $F\subset\R$ be finite. There is
$\delta>0$ such that $\Lm^1(\{x:x+sF\subseteq S\})>0$ for all $|s|<\delta$.
\end{lemma}

\begin{proof}
For $a\in F$ let $T_{a,s}=S-sa$. Translation is continuous in measure:
$\Lm^1(T_{a,s}\,\triangle\,S)\to0$ as $s\to0$, since $\|\mathbf 1_S(\cdot+h)-\mathbf 1_S\|_{L^1}\to0$
as $h\to0$. Choose $\delta>0$ with $\sum_{a\in F}\Lm^1(T_{a,s}\triangle S)<\Lm^1(S)$ for $|s|<\delta$.
Since
\[
  S\subseteq\{x: x+sF\subseteq S\}\ \cup\ \bigcup_{a\in F}\,(S\setminus T_{a,s}),
\]
we get $\Lm^1(S)\le\Lm^1(\{x:x+sF\subseteq S\})+\sum_{a\in F}\Lm^1(T_{a,s}\triangle S)$, and the
claim follows.
\end{proof}

\begin{proof}[Proof of Theorem~\ref{thm:A}]
$\Delta(E)$ is the image of the compact set $E\times E$ under the continuous map
$(x,y)\mapsto|x-y|$, so it is compact. In particular it is measurable with finite measure, and
$\Lm^1(\Delta(E))>0$ by Theorem~\ref{thm:inputFalconer}. Lemma~\ref{lem:steinhaus} with $S=\Delta(E)$ gives the first
claim. For the second, apply it with $s=\pm\min(\delta,\delta')/2$; a set of positive measure is
nonempty.
\end{proof}

\begin{remark}
The set $\{(x,s):x+sF\subseteq\Delta(E)\}$ is closed, and by Tonelli's theorem it has positive
planar measure. Theorem~\ref{thm:A} uses nothing about distance sets beyond compactness and positive measure. Its
point is the contrast with Theorem~\ref{thm:B} and Question~\ref{q:main}: the finite-to-infinite gap
of the Erd\H{o}s similarity problem reappears inside Falconer distance sets.
\end{remark}

\section{The pinned obstruction: proof of Theorem~\ref{thm:B}}

\begin{lemma}[Adding the limit point]\label{lem:zero}
Let $A\subset\R$ contain no nontrivial affine copy of $\Dy$. Then neither does $A\cup\{0\}$.
\end{lemma}

\begin{proof}
Suppose $x+s2^{-n}\in A\cup\{0\}$ for all $n\ge1$, with $s\ne0$. The points $x+s2^{-m}$, $m\ge0$,
are pairwise distinct. If none of them is $0$, then $x+s\Dy\subseteq A$, a contradiction. Otherwise
$x+s2^{-m}=0$ for exactly one $m\ge0$. The tail $\{x+s2^{-(m+n)}:n\ge1\}=x+(s2^{-m})\Dy$ is then
an affine copy of $\Dy$ avoiding $0$, hence contained in $A$, again a contradiction.
\end{proof}

\begin{lemma}[Radial sets]\label{lem:radial}
Let $d\ge1$ and let $A\subseteq[0,1]$ be compact with $\Lm^1(A)=c>0$. Then
$K_A=\{y\in\R^d:|y|\in A\}$ is compact, contained in the closed unit ball, and
$\Lm^d(K_A)\ge c^d\,\Lm^d(B(0,1))$.
\end{lemma}

\begin{proof}
Compactness and boundedness are clear. In polar coordinates,
\[
  \Lm^d(K_A)=d\,\Lm^d(B(0,1))\int_A r^{d-1}\,dr .
\]
Since $r\mapsto r^{d-1}$ is nondecreasing,
\begin{align*}
  \int_A r^{d-1}dr-\int_0^c r^{d-1}dr
  &=\int_{A\setminus[0,c]}r^{d-1}dr-\int_{[0,c]\setminus A}r^{d-1}dr\\
  &\ge c^{d-1}\bigl(\Lm^1(A\setminus[0,c])-\Lm^1([0,c]\setminus A)\bigr)=c^{d-1}(\Lm^1(A)-c)=0,
\end{align*}
so $\int_A r^{d-1}dr\ge c^d/d$.
\end{proof}

\begin{lemma}\label{lem:dim}
If $K\subset\R^d$ has $\Lm^d(K)>0$, then $\dimH K=d$.
\end{lemma}

\begin{proof}
On $\R^d$ the $d$-dimensional Hausdorff measure (normalized as in Mathlib, for the sup-norm) is a
positive multiple of $\Lm^d$, and the Euclidean and sup norms are bi-Lipschitz equivalent.
Hence $\mathcal H^d(K)>0$ and $\dimH K\ge d$. The reverse inequality is trivial.
\end{proof}

\begin{lemma}[Steinhaus]\label{lem:interval}
If $K\subset\R^d$ ($d\ge1$) is compact with $\Lm^d(K)>0$, then $[0,\varepsilon)\subseteq\Delta(K)$
for some $\varepsilon>0$. Consequently $\Delta(K)\supseteq\frac\varepsilon2\Dy$.
\end{lemma}

\begin{proof}
By Steinhaus' theorem $K-K\supseteq B(0,\varepsilon)$ for some $\varepsilon>0$. For
$0\le r<\varepsilon$ write $re_1=x-y$ with $x,y\in K$; then $r=|x-y|\in\Delta(K)$.
\end{proof}

\begin{proof}[Proof of Theorem~\ref{thm:B}]
Let $A$ be given by Theorem~\ref{thm:inputDyadic} for $\eta$, and put $K=\{y:|y|\in A\cup\{0\}\}$. The set
$A\cup\{0\}\subseteq[0,1]$ is compact, so $K$ is compact and lies in the closed unit ball, and
$0\in K$. Since $K\supseteq K_A$ and $\Lm^1(A)>1-\eta$, Lemma~\ref{lem:radial} gives
$\Lm^d(K)\ge(1-\eta)^d\Lm^d(B(0,1))>0$, and Lemma~\ref{lem:dim} gives $\dimH K=d$. Further,
$\Delta_0(K)=\{|y|:y\in K\}\subseteq A\cup\{0\}$, which contains no affine copy of $\Dy$ by
Lemma~\ref{lem:zero}. Lemma~\ref{lem:interval} shows that $\Delta(K)$ contains one.

For the last sentence, let $d\ge2$ and $p\ne0$, and pick $\rho\in A$ with $\rho>0$ (possible since
$\Lm^1(A)>0$). The sphere $\{|y|=\rho\}\subseteq K$ is connected and the continuous function
$y\mapsto|p-y|$ on it attains the values $|\rho-|p||$ and $\rho+|p|$. Hence $\Delta_p(K)$ contains
the interval $[\,|\rho-|p||,\rho+|p|\,]$ of length $2\min(\rho,|p|)>0$. Any nondegenerate interval
$[a,b]$ contains $a+(b-a)\Dy$.
\end{proof}

\begin{remark}[Other ratios]
The proof uses only two properties of $\Dy$: the avoidance property, and the fact that tails of $\Dy$
are affine copies of $\Dy$. Both hold for $\{q^n:n\ge1\}$ with $q\in(0,1)$. So Theorem~\ref{thm:B}
holds for every geometric sequence, conditional on the geometric case~\cite{OAIGeometric}, which is
not part of the Lean formalization.
\end{remark}

\section{Above the Mattila--Sj\"olin threshold: proof of Proposition~\ref{prop:C}}

If $\dimH E>(d+1)/2$, the Mattila--Sj\"olin theorem~\cite{MattilaSjolin1999} gives an interval
$[a,b]\subseteq\Delta(E)$ with $a<b$. If $\Lm^d(E)>0$, Lemma~\ref{lem:interval} gives one. For bounded
$P\subseteq[0,M]$ with $M>0$, we have $a+\frac{b-a}{M}P\subseteq[a,b]$. If $\Delta(E)$ contains no affine copy of
$\Dy$, it therefore has empty interior. If moreover $\dimH E>d/2$, then $\Lm^1(\Delta(E))>0$ by
Theorem~\ref{thm:inputFalconer}. \qed

\section{Open questions}\label{sec:open}
\begin{enumerate}[leftmargin=*]
\item Question~\ref{q:main} in the window $d/2<\dimH E\le(d+1)/2$. A positive answer would need
  multi-scale structure of distance sets beyond positive measure. A negative answer would need a
  construction controlling all distances of a set of dimension $>d/2$ simultaneously.
\item \emph{All pins.} Is there a compact $E\subset\R^d$ with $\Lm^d(E)>0$ (or merely $\dimH E>d/2$)
  such that $\Delta_p(E)$ contains no affine copy of $\Dy$ for \emph{every} $p\in E$? The complement
  of such a set must contain, around every point of $E$, entire spheres whose radii meet every
  affine copy of $\Dy$. This suggests a connection with spherical maximal function estimates.
  Pinned Falconer-type theorems~\cite{PeresSchlag2000,GIOW2020} give pins with large pinned distance sets,
  but say nothing about such patterns.
\item The same questions for general infinite sets $P$, for instance for every $P$ known to be
  non-universal.
\end{enumerate}

\section{Formalization}\label{sec:lean}
\begin{sloppypar}

The Lean development accompanies this note. It is pinned to Lean~4 \texttt{v4.34.1} and
Mathlib commit \texttt{d13f23b}, the same versions as OpenAI's library, so the input statements
typecheck identically. It consists of three files:
\begin{itemize}[leftmargin=*]
\item \texttt{Statements.lean}: the propositions \texttt{FalconerStatement} and
  \texttt{DyadicAvoidanceStatement}, copied verbatim from OpenAI's statement files
  \texttt{FalconerAllDimensions.lean} and \texttt{DyadicAvoidance.lean}, and the definitions of distance
  sets and of ``contains an affine copy of $\Dy$''.
\item \texttt{FinitePatterns.lean}: Lemma~\ref{lem:steinhaus} and Theorem~\ref{thm:A}
  (\texttt{finite\_pattern\_distSet}, \texttt{exists\_finite\_pattern\_distSet}).
\item \texttt{PinnedObstruction.lean}: Lemmas~\ref{lem:zero}, \ref{lem:dim}, \ref{lem:interval},
  positivity of the volume in Lemma~\ref{lem:radial}, and Theorem~\ref{thm:B}
  (\texttt{pinned\_dyadic\_obstruction}). The formal statement is:
\end{itemize}
\begin{alltt}\footnotesize
theorem pinned_dyadic_obstruction (hDy : DyadicAvoidanceStatement) \{d : \(\mathbb{N}\)\}
    (hd : 1 \(\le\) d) \{\(\eta\) : \(\mathbb{R}\)\} (h\(\eta\)0 : 0 < \(\eta\)) (h\(\eta\)1 : \(\eta\) < 1) :
    \(\exists\) K : Set (EuclideanSpace \(\mathbb{R}\) (Fin d)),
      IsCompact K \(\wedge\) K \(\subseteq\) closedBall 0 1 \(\wedge\)
      (0 : EuclideanSpace \(\mathbb{R}\) (Fin d)) \(\in\) K \(\wedge\) 0 < volume K \(\wedge\) dimH K = d \(\wedge\)
      \(\neg\) ContainsDyadicCopy (pinnedDistSet 0 K) \(\wedge\) ContainsDyadicCopy (distSet K)
\end{alltt}
Not formalized are the explicit constant $(1-\eta)^d$ in Lemma~\ref{lem:radial} (only positivity
of the volume), the last sentence of Theorem~\ref{thm:B}, the remark on other ratios, and
Proposition~\ref{prop:C}, which relies on the Mattila--Sj\"olin theorem.

\emph{Re-checking the inputs.} A script in the repository copies, from a checkout of
\texttt{github.com/openai/math} at commit \texttt{adc7f12}, the transitive closure of OpenAI's Lean
modules proving the two inputs: 2459 files, depending only on Mathlib. These are built with the
same toolchain. A further module then states
\texttt{theorem : FalconerStatement := OAI.Falconer.falconer\_distance\_conjecture} and the
analogous dyadic statement, and derives unconditional versions of Theorems~\ref{thm:A} and~\ref{thm:B}.
The kernel accepts these, and \texttt{\#print axioms} reports only \texttt{propext},
\texttt{Classical.choice} and \texttt{Quot.sound}. We additionally replayed the resulting environment
with Lean's \texttt{leanchecker}. We did not run OpenAI's \texttt{comparator} setup. A formal proof
certifies the statement as formalized. Readers should still compare the formal statements with
the informal ones.
\end{sloppypar}

\section*{Acknowledgments and use of AI}
This note builds on results produced by an AI model and released by OpenAI~\cite{OAIFalconer,OAIDyadic}.
It was itself prepared with substantial assistance from Claude (Anthropic), an AI system, which
helped choose the combination studied here, search the literature, draft the text, and write
the Lean formalization. The author checked the mathematics and the formal statements and takes
full responsibility for the content.

\bibliographystyle{amsplain}
\bibliography{refs}

\end{document}
